problem:  
Let \( (M^{2n}, g) \) be a compact, simply connected Riemannian manifold with **nonnegative sectional curvature**, supporting an effective, isometric, **slice‑maximal torus action** by \(T^k\).  
Let \( (\mathcal{O}, d_g, \lambda) = (M/T^k, d_g, \lambda) \) denote the corresponding **labeled Alexandrov orbifold**, where \(\lambda\) records the primitive isotropy weights along the codimension‑two strata.

Fix the action and its labeling data \(\lambda\), and consider the space  
\[
\mathcal{M}_{\ge 0}^{T^k}(M) = \{ T^k\text{-invariant Riemannian metrics } g \text{ on } M \mid \sec_g \ge 0 \}.
\]
For each \(g\), the orbit space acquires an Alexandrov metric \(d_g\) of curvature \(\ge 0\) on \(\mathcal{O}\). Equip the set  
\[
\mathfrak{Q}_{\ge 0}(\mathcal{O}, \lambda)
\]
of such metrics \(d_g\) with the **Gromov–Hausdorff topology**, where isometries are required to respect the labeled stratification of \(\mathcal{O}\).

Define the **quotient map**
\[
\Pi: \mathcal{M}_{\ge 0}^{T^k}(M) \longrightarrow \mathfrak{Q}_{\ge 0}(\mathcal{O}, \lambda), \qquad g \mapsto d_g.
\]

**Open problem (Local Gromov–Hausdorff rigidity of Alexandrov quotient metrics):**  
Does there exist a metric \(g_0 \in \mathcal{M}_{\ge0}^{T^k}(M)\) for which \(\Pi\) is **locally constant** in a neighborhood of \(g_0\)?  
Equivalently, are there nonnegatively curved, \(T^k\)-invariant metrics such that any sufficiently small curvature‑preserving, \(T^k\)-invariant perturbation of the metric produces an **isometric Alexandrov quotient** \((\mathcal{O}, d_g)\) (up to labeled isometry)?  
If not, can one describe the directions in which infinitesimal perturbations of \(g\) change the quotient metric structure while keeping slice‑maximality?

Why is it a "good" problem:  
This formulation avoids undefined analytic structures but keeps the core geometric idea intact: the search for **metric rigidity of Alexandrov quotients** under curvature‑preserving invariant deformations.

* **Mathematical soundness:** The spaces involved—\(T^k\)-invariant nonnegatively curved metrics on \(M\) and their Alexandrov quotients with Gromov–Hausdorff topology—are standard objects. The notion of local constancy of \(\Pi\) is well‑defined and meaningful without requiring a differentiable moduli structure.

* **Conceptual bridge:** It directly connects
  - deformation theory of invariant Riemannian metrics under curvature constraints,
  - Alexandrov geometry of quotient metric spaces,
  - classification theory of slice‑maximal torus actions.
  The question explores how the **large‑scale metric geometry of the quotient** reacts to **fine‑scale curvature deformations** upstairs.

* **Novelty and depth:** The classification of slice‑maximal torus actions describes \(M\) combinatorially, but **the metric geometry of the quotient remains uncharted**. Whether it is locally rigid or flexible under curvature‑preserving perturbations is unknown, and its resolution would broaden understanding of both nonnegative curvature moduli and Alexandrov rigidity.

* **Analogy and potential impact:** A positive answer would parallel Mostow‑type rigidity, but now for **positively curved quotient geometries**; a negative answer would reveal hidden **metric moduli phenomena** invisible topologically.

* **Simplicity:** The question can be restated succinctly:  
> “Can the Alexandrov metric on the quotient of a nonnegatively curved, slice‑maximal torus action deform under curvature‑preserving invariant perturbations, or is it locally rigid?”  
This simple yet profound question sits at the intersection of Riemannian, Alexandrov, and transformation‑group geometry, making it a paradigmatic “good” research‑level problem.